\section{Model and theoretical framework}
\label{sec:theory}

The interacting exciton and biexciton ground-state energies are obtained with
the variational method: trial correlation factors are optimized by minimizing
the corresponding Hamiltonian expectation values. We first establish the
single-particle confinement states on which these trial states are built.

\subsection{Geometry, units, and single-particle ground state}
\label{sec:single-particle}

We consider a hard-wall semi-ellipsoidal quantum dot with cylindrical
coordinates $\vect r=(r,\phi,z)$ and confinement potential
\begin{equation}
 V_{\mathrm{conf}}(r,z)=
 \begin{cases}
  0, & r^2/a^2+z^2/c^2\leq 1\ \text{and}\ z\geq0,\\
  \infty, & \text{otherwise},
 \end{cases}
 \qquad c/a\ll1 .
 \label{eq:confinement-potential}
\end{equation}
The in-plane and axial semiaxes are $a$ and $c$, respectively, as illustrated
in \cref{fig:semi-ellipsoid}.

\begin{figure}[ht]
 \centering
 \includegraphics[width=0.52\linewidth]{Semi-ellipsoid.jpeg}
 \caption{Semi-ellipsoidal quantum-dot geometry, with lateral semiaxis $a$
 and axial semiaxis $c$.}
 \label{fig:semi-ellipsoid}
\end{figure}

Lengths and energies are expressed in the effective electron Bohr radius and
Rydberg,
\begin{equation}
 \rB=\frac{4\pi\epsilon_0\epsilon_r^{(E)}\hbar^2}{\meff e^2},
 \qquad
 E_{\mathrm R}=\frac{\hbar^2}{2\meff\rB^2}.
 \label{eq:effective-units}
\end{equation}
For the GaAs energy calculation we use the standard bulk parameter set
\cite{vurgaftman2001},
$\epsilon_r^{(E)}=12.8$, $\meff=0.067m_0$, and
$\mheff=0.45m_0$, giving
$\rB=\SI{10.109654}{nm}$ and
$E_{\mathrm R}=\SI{5.563851565}{meV}$.

Because the axial confinement is much stronger than the in-plane
confinement, the single-particle problem is treated adiabatically
\cite{hayrapetyan2012}. At fixed $r$, the available height is
\begin{equation}
 z_{\max}(r)=c\sqrt{1-r^2/a^2},
 \label{eq:zmax}
\end{equation}
and the lowest axial factor is
\begin{equation}
 \chi_1(z;r)=
 \sqrt{\frac{2}{z_{\max}(r)}}
 \sin\left[\frac{\pi z}{z_{\max}(r)}\right].
 \label{eq:axial-wavefunction}
\end{equation}
For a carrier of mass $m_q$ with $q=e,h$ the axial eigenenergy acts as the adiabatic
potential for the slow planar motion,
\begin{equation}
 U_q(r)=\frac{\meff}{m_q}\frac{\pi^2}{z_{\max}^2(r)}
 =\frac{\meff}{m_q}\frac{\pi^2}{c^2}
   \left(1+\frac{r^2}{a^2}+O(r^4/a^4)\right),
 \label{eq:adiabatic-potential}
\end{equation}
where lengths and energies are in the effective units of
\cref{eq:effective-units}. Thus, after subtracting the constant axial energy,
the planar harmonic-oscillator Hamiltonian is
\begin{equation}
 H_{q,\parallel}=\frac{\meff}{m_q}
 \left[-\nabla_{\parallel}^2+\frac{\pi^2r^2}{a^2c^2}\right].
 \label{eq:planar-hamiltonian}
\end{equation}
Writing $\lambda=\pi/(ac)$, its normalized ground-state component and the
corresponding axial and planar ground-state energies are
\begin{equation}
 \varphi_{q,\parallel}(r)=\sqrt{\frac{\lambda}{\pi}}
 e^{-\lambda r^2/2},\qquad
 \varepsilon_{q,z}=\frac{\meff}{m_q}\frac{\pi^2}{c^2},\qquad
 \varepsilon_{q,\parallel}=\frac{\meff}{m_q}\frac{2\pi}{ac}.
 \label{eq:single-particle-scales}
\end{equation}
The total ground-state confinement energy of carrier $q$ is therefore
\begin{equation}
 E_q=\varepsilon_{q,z}+\varepsilon_{q,\parallel}
 =\frac{\meff}{m_q}\left(\frac{\pi^2}{c^2}+\frac{2\pi}{ac}\right).
 \label{eq:carrier-ground-energy}
\end{equation}
The normalized carrier ground orbital is therefore
$\psi_q(r,\phi,z)=\chi_1(z;r)\varphi_{q,\parallel}(r)$. Thus the uncorrelated exciton and biexciton confinement energies are
\begin{equation}
 E_X^{(0)}=E_e+E_h,\qquad
 E_{XX}^{(0)}=2E_e+2E_h=2E_X^{(0)}.
 \label{eq:single-particle-energies}
\end{equation}
The band gap is not included in these energies.

\subsection{Correlated exciton}
\label{sec:exciton-model}

We denote electron coordinates by $1,2$ and hole coordinates by $a,b$; the
ground exciton uses coordinates $1$ and $a$. Pair distances are the
three-dimensional Euclidean distances $r_{ij}=|\vect r_i-\vect r_j|$. In the
effective units of \cref{eq:effective-units}, its Hamiltonian is
\begin{equation}
 H_X=h_e(1)+h_h(a)-\frac{2}{r_{1a}},
 \qquad
 h_e\psi_e=E_e\psi_e,\quad h_h\psi_h=E_h\psi_h.
 \label{eq:exciton-hamiltonian}
\end{equation}
The uncorrelated state $\Psi_{0X}=\psi_e(1)\psi_h(a)$ is therefore an
eigenstate of $h_e+h_h$ with energy $E_X^{(0)}=E_e+E_h$. Electron--hole
correlation is introduced through
\begin{equation}
 \Psi_X=\Psi_{0X}F_X,\qquad F_X=\exp(-\alpha_X r_{1a}),
 \label{eq:exciton-trial}
\end{equation}
where $\alpha_X>0$ is varied independently at each geometry. The variational
energy is evaluated from
\begin{equation}
 E_X=\frac{\langle\Psi_X|H_X|\Psi_X\rangle}
 {\langle\Psi_X|\Psi_X\rangle}.
 \label{eq:exciton-expectation}
\end{equation}
The denominator is
\begin{equation}
 \NX=\int
 |\Psi_{0X}|^2F_X^2\dd\tau
 =\int|\Psi_{0X}|^2e^{-2\alpha_X r_{1a}}\dd\tau .
 \label{eq:exciton-norm}
\end{equation}
Using the single-particle eigenvalue equations in
\cref{eq:exciton-hamiltonian} and integrating the kinetic terms by parts gives
\begin{equation}
 E_X=E_X^{(0)}+\frac{1}{\NX}\int|\Psi_{0X}|^2
 \left[|\nabla_1F_X|^2+\frac{\meff}{\mheff}|\nabla_aF_X|^2
 -\frac{2F_X^2}{r_{1a}}\right]\dd\tau .
 \label{eq:exciton-energy-integral}
\end{equation}
Since $|\nabla_1r_{1a}|=|\nabla_ar_{1a}|=1$, the kinetic contribution
generated by $F_X$ is
\begin{equation}
 K_X(\alpha_X)=
 \alpha_X^2\left(1+\frac{\meff}{\mheff}\right).
 \label{eq:exciton-kinetic}
\end{equation}
Defining the Coulomb numerator as
\begin{equation}
 V_X=\int
 |\Psi_{0X}|^2F_X^2
 \left(-\frac{2}{r_{1a}}\right)\dd\tau
 \label{eq:exciton-coulomb}
\end{equation}
reduces \cref{eq:exciton-energy-integral} to
\begin{equation}
 \EX=E_X^{(0)}+\Delta E_X,\qquad
 \Delta E_X=K_X+\frac{V_X}{\NX}.
 \label{eq:exciton-energy}
\end{equation}
For each dot geometry, the remaining integrals $N_X$ and $V_X$ are evaluated
numerically.

\subsection{Correlated biexciton}
\label{sec:biexciton-model}

The two-electron--two-hole interaction in effective units is
\begin{equation}
 V_{XX}=2\left(
 \frac{1}{r_{12}}+\frac{1}{r_{ab}}
 -\frac{1}{r_{1a}}-\frac{1}{r_{1b}}
 -\frac{1}{r_{2a}}-\frac{1}{r_{2b}}
 \right).
 \label{eq:biexciton-coulomb}
\end{equation}
For the biexciton we use the ansatz~\cite{}
\begin{equation}
 \Psi_{XX}=\Psi_{0XX}F_{XX},\qquad
 \Psi_{0XX}=\psi_e(1)\psi_e(2)\psi_h(a)\psi_h(b),
 \label{eq:biexciton-trial}
\end{equation}
with the exchange-symmetric correlation factor
\begin{equation}
\begin{aligned}
 F_{XX}&=G(P+Q),\\
 G&=r_{ab}^{\gamma}e^{-\delta r_{ab}},\\
 P&=\exp\left[-\alpha(r_{1a}+r_{2b})
                 -\beta(r_{1b}+r_{2a})\right],\\
 Q&=\exp\left[-\beta(r_{1a}+r_{2b})
                 -\alpha(r_{1b}+r_{2a})\right].
\end{aligned}
 \label{eq:correlation-factor}
\end{equation}
The four variational parameters $(\alpha,\beta,\gamma,\delta)$ control the
direct and cross electron--hole pairings and the hole--hole correlation.

To write the kinetic numerator compactly, define
$\vect g_i^P=\nabla_i\ln(GP)$ and
$\vect g_i^Q=\nabla_i\ln(GQ)$. For electron coordinates the derivatives of
$G$ vanish, whereas for hole coordinates they contain the logarithmic
derivative of $G$ in \cref{eq:correlation-factor}. Unlike the exciton kinetic contribution,
the biexciton kinetic numerator has no analytical form and is written as
\begin{equation}
 K_{XX}=\int|\Psi_{0XX}|^2G^2\left[\sum_{i=1,2}|P\vect g_i^P+Q\vect g_i^Q|^2+\frac{\meff}{\mheff}\sum_{j=a,b}|P\vect g_j^P+Q\vect g_j^Q|^2\right]\dd\tau .
 \label{eq:biexciton-kinetic-numerator}
\end{equation}
The energy is then
\begin{equation}
 \EXX=E_{XX}^{(0)}+\Delta E_{XX},\qquad
 \Delta E_{XX}=
 \frac{K_{XX}+\langle\Psi_{XX}|V_{XX}|\Psi_{XX}\rangle}
 {\langle\Psi_{XX}|\Psi_{XX}\rangle}.
 \label{eq:biexciton-energy}
\end{equation}
In the numerical calculation, the kinetic contribution in
\cref{eq:biexciton-kinetic-numerator} and the potential contribution are
combined and evaluated in a single integral. No direct/cross or
particle-exchange multiplicities are introduced in the final estimator.

We define the biexciton binding energy as
\begin{equation}
 \Ebind=2\EX-\EXX
       =2\Delta E_X-\Delta E_{XX}.
 \label{eq:binding-energy}
\end{equation}
With this convention $\Ebind>0$ denotes a bound biexciton. The
single-particle part cancels exactly, making \cref{eq:binding-energy} a small
difference of correlation corrections.

\subsection{Radiative lifetimes}
\label{sec:observables}

Following Hayrapetyan \textit{et al.}\ \cite{hayrapetyan2019}, we calculate the exciton and biexciton radiative lifetimes from
\begin{equation}
 \tauX=
 \frac{2\pi\epsilon_0\meff c_0^3\hbar^2}
 {\sqrt{\epsilon_r^{(\mathrm{opt})}}\,e^2
 E_{\mathrm{exc}}^2f_X},
 \qquad
 \tauXX=\frac{\tauX}{4}.
 \label{eq:lifetimes}
\end{equation}
The relation $\tauXX=\tauX/4$ follows the biexciton decay-channel convention
of Ref.~\cite{hayrapetyan2019}; it is not an independently optimized quantity.

Evaluation of $\tauX$ requires the oscillator strength and transition energy,
which are
\begin{equation}
 f_X=\frac{E_P}{E_{\mathrm{exc}}}|\MX|^2,\qquad
 E_{\mathrm{exc}}=E_g+\EX,
 \label{eq:oscillator-strength}
\end{equation}
where, $E_g$ is the band gap. The remaining quantity $\MX$ is the electron--hole coincidence overlap. At
coincidence the exciton Jastrow factor in \cref{eq:exciton-trial} is unity,
and because the electron and hole ground orbitals are matched in the present
model, the normalized overlap reduces to
\begin{equation}
 \MX=\frac{1}{\sqrt{\NX}},\qquad
 |\MX|^2=\frac{1}{\NX}.
 \label{eq:coincidence-overlap}
\end{equation}
Thus the lifetime can be obtained from the already calculated exciton norm;
no additional oscillator-strength integral is required.

For the optical calculation we take
$E_g=\SI{1.5192}{eV}$, $\meff=0.0665m_0$,
$E_P=\SI{22.71}{eV}$, and
$\epsilon_r^{(\mathrm{opt})}=12.91$, matching the parameter set of
Ref.~\cite{hayrapetyan2019}. We keep this optical set distinct from the bulk
parameters in \cref{eq:effective-units} because the latter define the units
and single-particle states used in the variational energy calculation,
whereas the former enter the dimensional radiative-lifetime parametrization.
This separation preserves the energy calculation while reproducing the
cited optical calibration.
